\section{Function-space topology from explicit finite-rank operators}
\label{sec:function-topology}

The earlier topology packet already supplies covers, carriers, finite
homology calculations, and a careful warning about coherent witnesses.
The next useful step is to construct \emph{continuous functions} from
finite observations, with all denominators and continuity constants
controlled.  We give such a construction on an arbitrary presented
totally bounded metric space.  It yields a quantitative form of
Arzel\`a--Ascoli, a converse, and a gluing theorem.  These results concern
infinite spaces and function spaces even though their arithmetic proofs
use finite lists.

\subsection{The input presentation and the meaning of the statement}

Write $d_s=2^{-s}$.  A rational metric basis is a decidable collection
$A\subseteq\N$ with a fixed PR rational-valued distance $d(a,b)$ satisfying
the pseudometric laws.  Zero distance is quotiented extensionally.
For each $s$, the presentation supplies a positive natural number $N(s)$
and a list
\[
 c(s,0),\ldots,c(s,N(s)-1)\in A
\]
with the covering assertion
\begin{equation}\label{eq:top-net}
 a\in A\quad\Longrightarrow\quad
 \exists i<N(s)\;d(a,c(s,i))\le d_s.
\end{equation}
The list functions and $N$ are fixed PR programs, and the covering
assertion is an explicit premise or a previously derived theorem.
Let $X$ denote the metric completion in the external interpretation.
The list at scale $s+1$ covers $X$ at radius $d_s$ constructively:
given a regular name for $x$, take a rational approximation $a$ within
$d_{s+1}$ of $x$, and then use \eqref{eq:top-net} to find a list
member within $d_{s+1}$ of $a$.  The triangle inequality gives the
claimed bound.  This extra precision bit avoids selecting a minimum
among finitely many real distances.  All arithmetic calculations
below are performed at rational queries; extension to $X$ uses their
displayed continuity moduli.

For readers wishing to avoid any use of a real-valued supremum, every
assertion $\|f-g\|_\infty\le\eta$ below means
$\forall x\in X\;|f(x)-g(x)|\le\eta$.  It does not invoke existence of a
supremum as a new mathematical principle.

Fix a family $(f_a)_{a\in D}$ of real-valued functions on $X$, an integer
$B\ge0$, and a fixed PR modulus $\mu$ such that
\begin{equation}\label{eq:top-family}
 |f_a(x)|\le B,\qquad
 d(x,y)\le 2^{-\mu(j)}\quad\Longrightarrow\quad
 |f_a(x)-f_a(y)|\le 2^{-j}.
\end{equation}
The same modulus must work for every family member.  The symbol $a$ is
a parameter for one fixed evaluator, not a code fed to a universal PR
interpreter.  At the arithmetic level, observations of $f_a$ are
rationals $E(a,x,p)$ with error at most $d_p$ at a rational query
$x\in A$.  A fully rational specification is given below.  The family
need not have decidable equality, and its members need not be specified
by a single finite polynomial or by a discrete table.

\subsection{A denominator bound produced by the net}

\begin{lemma}[Quantitative partition of unity]\label{lem:tent-partition}
Let $r>0$ be rational and let $c_0,\ldots,c_{n-1}$ be an $r/2$-net in
$X$, with $n\ge1$.  Define
\[
 w_i(x)=\max\{0,r-d(x,c_i)\},\qquad
 W(x)=\sum_{i<n}w_i(x),\qquad
 \lambda_i(x)=\frac{w_i(x)}{W(x)}.
\]
Then $W(x)\ge r/2$, $\lambda_i(x)\ge0$, and
$\sum_{i<n}\lambda_i(x)=1$.  Moreover,
\begin{equation}\label{eq:tent-lipschitz}
 \sum_{i<n}|\lambda_i(x)-\lambda_i(y)|
 \le\frac{4n}{r}d(x,y).
\end{equation}
If $h=(h_i)_{i<n}$ and $|h_i|\le B$, the reconstructed function
\[
 R(h)(x)=\sum_{i<n}\lambda_i(x)h_i
\]
is bounded by $B$ and is Lipschitz with constant $4Bn/r$.
For two coefficient vectors,
\begin{equation}\label{eq:tent-coefficient}
 |R(h)(x)-R(\widetilde h)(x)|
 \le\max_{i<n}|h_i-\widetilde h_i|.
\end{equation}
\end{lemma}

\begin{proof}
Choose a covering index $j$ for $x$.  Then
$w_j(x)=r-d(x,c_j)\ge r/2$, and all other weights are nonnegative.
This proves the denominator bound and the first three assertions.
The triangle inequality gives
$|d(x,c_i)-d(y,c_i)|\le d(x,y)$.  The real function
$t\mapsto\max(0,r-t)$ does not increase distances, so
\[
 |w_i(x)-w_i(y)|\le d(x,y),\qquad
 |W(x)-W(y)|\le n\,d(x,y).
\]
Using positive denominators, add and subtract $w_i(y)/W(x)$ to obtain
\[
 |\lambda_i(x)-\lambda_i(y)|
 \le\frac{|w_i(x)-w_i(y)|}{W(x)}
  +\frac{w_i(y)|W(y)-W(x)|}{W(x)W(y)}.
\]
Sum this inequality.  In the second term the sum of $w_i(y)$ is
$W(y)$, and hence
\[
 \sum_i|\lambda_i(x)-\lambda_i(y)|
 \le\frac{n\,d(x,y)+|W(y)-W(x)|}{W(x)}
 \le\frac{4n}{r}d(x,y).
\]
Multiply by the coefficient bound $B$ to obtain the Lipschitz assertion.
The bound $|R(h)(x)|\le B$ and \eqref{eq:tent-coefficient} follow by
nonnegativity of the weights and their sum being one.  No cancellation,
choice of a nearest centre, or decision about equality of real points
is used.
\end{proof}

The formula is a continuous partition of unity; smoothness is not
claimed.  On the rational basis, all weights and all denominators are
exact rational numbers.  The inverse operation is justified by a
\emph{derived} lower bound $r/2$, rather than by a new oracle promising
that a denominator is nonzero.

\subsection{An explicit quantitative Arzel\`a--Ascoli theorem}

\begin{theorem}[Finite continuous approximants with an entropy bound]
\label{thm:quant-ascoli}
Under \eqref{eq:top-net} and \eqref{eq:top-family}, put
\begin{align*}
 s_k&=\mu(k+2)+2,& r_k&=2^{-\mu(k+2)},& n_k&=N(s_k),\\
 \delta_k&=2^{-(k+3)},&
 M_k&=2B\,2^{k+3}+1,& H_k&=M_k^{n_k}.
\end{align*}
The basis list has mesh $r_k/4$ and hence covers $X$ with mesh $r_k/2$.
Using the centres $c(s_k,i)$, form the weights of
Lemma~\ref{lem:tent-partition} with radius $r_k$.
There is an explicit list of $H_k$ continuous functions $g_{k,h}$
such that
\begin{enumerate}
 \item $|g_{k,h}(x)|\le B$ and each $g_{k,h}$ is Lipschitz with
 constant at most $4Bn_k/r_k$;
 \item for every $a\in D$, finitely many evaluations of $f_a$ at the
 $n_k$ net points, each to accuracy $\delta_k$, produce an index
 $h<H_k$ satisfying
 \begin{equation}\label{eq:ascoli-cover}
  \|f_a-g_{k,h}\|_\infty\le 2^{-(k+1)};
 \end{equation}
 \item the list size, its rational evaluators, the selected index, and
 every bound above are PR in $k,B$ and the numerical presentation data,
 for each fixed collection of input programs.
\end{enumerate}
The centres of this cover lie in the ambient space of continuous
functions.  They are not asserted to belong to the original family.
\end{theorem}

\begin{proof}
The grid
\[
 G_k=\{-B,-B+\delta_k,\ldots,B-\delta_k,B\}
\]
has exactly $M_k$ members because $B$ is an integer and
$\delta_k^{-1}=2^{k+3}$.  If $B=0$, the single vector is the zero
vector.  Otherwise enumerate all vectors in $G_k^{n_k}$ by
base-$M_k$ digits.  For such a vector $h$, define
$g_{k,h}=R(h)$.  Lemma~\ref{lem:tent-partition} gives the first
assertion.  Each list element has a rational evaluator on $A$ given by
the displayed finite sum and quotient; the Lipschitz bound extends it
uniquely to the completion.

For a particular $f_a$, obtain a rational $e_i$ satisfying
$|e_i-f_a(c(s_k,i))|\le\delta_k$.  Define
\[
 h_i=\operatorname{clip}_{[-B,B]}
       \bigl(\delta_k\lfloor e_i/\delta_k\rfloor\bigr).
\]
This is a decidable rational operation, including at grid boundaries.
Projection onto an interval does not increase distance from a point of
that interval.  Rounding down changes a rational by at most $\delta_k$.
Since $f_a(c(s_k,i))\in[-B,B]$, it follows that
\begin{equation}\label{eq:ascoli-sample-error}
 |h_i-f_a(c(s_k,i))|\le 2\delta_k=2^{-(k+2)}.
\end{equation}
The resulting vector lies in $G_k^{n_k}$, so its base-$M_k$ code is
less than $H_k$.

First fix $x\in A$, so that every support comparison is decidable.
If $w_i(x)>0$, then
$d(x,c(s_k,i))<r_k=2^{-\mu(k+2)}$.  The common modulus therefore gives
\[
 |f_a(c(s_k,i))-f_a(x)|\le 2^{-(k+2)}.
\]
For these indices, \eqref{eq:ascoli-sample-error} implies
$|h_i-f_a(x)|\le 2^{-(k+1)}$.  Indices with zero weight contribute
zero.  Taking the weighted average yields
\[
 |g_{k,h}(x)-f_a(x)|
 \le\sum_i\lambda_i(x)|h_i-f_a(x)|
 \le 2^{-(k+1)}\sum_i\lambda_i(x)=2^{-(k+1)}.
\]
This estimate holds for every rational query.  For a regular name of
an arbitrary $x\in X$, the supplied modulus for $f_a$ and the proved
Lipschitz bound for $g_{k,h}$ pass the same non-strict estimate to $x$.
No sign decision about a real distance is needed.

For the PR claim, exponentiation, finite list formation, base expansion,
signed rational floors, clipping, and bounded sums are PR.  The call
depth of the input evaluator is fixed; only its numerical arguments
vary.  Every loop is bounded by a displayed natural number, although
$H_k$ can be very large.  In particular no test of the form ``does this
function agree with some earlier function?'' is part of the enumeration.
\end{proof}

\begin{remark}[What the theorem supplies]
This is a uniform finite cover in the supremum-distance sense.  If one
uses a definition of total boundedness demanding that the centres
belong to the given family, additional witnesses for the inhabited
profile classes are needed.  Given one member of the family in each
inhabited profile class, the triangle inequality changes the radius
from $2^{-(k+1)}$ to $2^{-k}$.  The displayed construction itself does
not decide which profile classes are inhabited.
Likewise, replacing the supplied common modulus by bare pointwise
equicontinuity is an extra theorem; it is not a consequence of the
finite enumeration.
\end{remark}

\begin{corollary}[Finite-rank approximation of continuous functions]
\label{cor:finite-rank}
With the same weights, let
\[
 (P_k f)(x)=\sum_{i<n_k}\lambda_i(x)f(c(s_k,i)).
\]
Then $P_k$ is a positive linear operator, preserves constants, has
operator norm at most one, and its image is contained in the span of
$\lambda_0,\ldots,\lambda_{n_k-1}$.  Every function satisfying the
given modulus obeys
\[
 \|P_k f-f\|_\infty\le 2^{-(k+2)}.
\]
Thus the metric presentation produces finite-rank approximation
operators with an explicit convergence rate on every family with a
common modulus.
\end{corollary}

\begin{proof}
Linearity is distributivity of a finite sum.  Positivity, preservation
of constants, and the norm bound follow from the nonnegative weights
adding to one.  The range statement is immediate from the formula.
At a rational query $x$, every nonzero weight comes from a centre within
$2^{-\mu(k+2)}$ of $x$.  Apply the modulus to each such summand and
average, exactly as in the preceding proof but without quantization;
then extend the estimate along regular names.
\end{proof}

The operator statement uses ordinary continuous-function language in
the external interpretation.  Its arithmetic content is the family
of finite sum identities and error estimates.  It does not assert that
all continuous functions form a numerical sort of \ipra.

For example, on $[0,1]^d$ with the maximum metric, the dyadic grid
gives $N(s)=(2^s+1)^d$.  Functions bounded by the integer $B$ and
Lipschitz with constant at most $2^\ell$ have
$\mu(j)=j+\ell$.  The theorem therefore gives the entirely explicit
covering bound
\[
 H_k=(2B\,2^{k+3}+1)^{(2^{k+\ell+4}+1)^d}.
\]
The estimate is deliberately conservative.  Its purpose is to expose
the full primitive-recursive dependence and provide a reusable proof,
not to claim optimal metric entropy.

\subsection{The converse: finite approximants recover regularity}

\begin{proposition}[A quantitative converse]\label{prop:ascoli-converse}
Suppose that for every $j$ a finite nonempty list
$g_{j,0},\ldots,g_{j,L_j-1}$ covers a family $\mathcal F$ within
supremum distance $2^{-j}$.  Suppose each $g_{j,i}$ has an explicit
uniform modulus $\nu(j,i,t)$, and the list at $j=0$ has integer bounds
$|g_{0,i}|\le C_i$.  Then all $f\in\mathcal F$ satisfy
\[
 |f|\le1+\max_{i<L_0}C_i
\]
and share the modulus
\begin{equation}\label{eq:ascoli-converse-modulus}
 \mu_*(k)=\max_{i<L_{k+2}}\nu(k+2,i,k+1).
\end{equation}
If the displayed data are given by fixed PR programs, these new data
are PR as well.  It is enough to have the finite-cover assertion;
an algorithm choosing a covering index is not required for this
direction.
\end{proposition}

\begin{proof}
For the range bound, take a covering index for $f$ in the list at
$j=0$ and use the triangle inequality.  For the modulus, take a
covering index $i$ at $j=k+2$.  If
$d(x,y)\le2^{-\mu_*(k)}$, it is at most
$2^{-\nu(k+2,i,k+1)}$.  Consequently
\begin{align*}
 |f(x)-f(y)|
 &\le |f(x)-g_{k+2,i}(x)|
   +|g_{k+2,i}(x)-g_{k+2,i}(y)|
   +|g_{k+2,i}(y)-f(y)|\\
 &\le 2^{-(k+2)}+2^{-(k+1)}+2^{-(k+2)}=2^{-k}.
\end{align*}
The index is eliminated from a finite existential hypothesis after
this estimate; no countable choice is used.  Maxima of finite lists
are computed by a bounded recursion.
\end{proof}

The forward and reverse results identify the effective content of
Ascoli approximation.  Uniform range and continuity information
produce finite approximants, and finite approximants with their own
regularity information recover a common modulus.  Convergence of an
arbitrary sequence is a separate question.

\subsection{Gluing guarded local data}

\begin{theorem}[Continuous gluing with a produced modulus]
\label{thm:guarded-gluing}
Let $c_0,\ldots,c_{n-1}$ be an $r/2$-net.  For each $i$, suppose a
local function $f_i$ is given on the open ball
$\{x:d(x,c_i)<2r\}$.  Assume that on this ball $|f_i|\le B$ and
$f_i$ is Lipschitz with constant $L$.  Assume also that on every
overlap where both distances are less than $r$,
\[
 |f_i(x)-f_j(x)|\le\alpha.
\]
On the rational basis define $p_i(x)=w_i(x)f_i(x)$ when
$d(x,c_i)<r$ and $p_i(x)=0$ otherwise; these are decidable rational
support tests.  Here $w_i$ is as in Lemma~\ref{lem:tent-partition}.
The products have unique Lipschitz extensions to $X$.  For those
extensions,
\[
 g(x)=\frac{\sum_{i<n}p_i(x)}{W(x)}
\]
is a globally defined function with $|g|\le B$,
\begin{equation}\label{eq:gluing-lip}
 |g(x)-g(y)|\le
 \left(2nL+\frac{4nB}{r}\right)d(x,y),
\end{equation}
and $|g(x)-f_j(x)|\le\alpha$ whenever $d(x,c_j)<r$.
In particular exact overlap agreement yields a genuine global
continuous gluing of the restrictions to the $r$-balls.  The assertions
also hold for complex values, with
the absolute value replaced by the complex norm.
\end{theorem}

\begin{proof}
First take $x,y$ in the rational basis.  The function $p_i$ is
well-defined because the support where its
formula uses $f_i$ lies strictly inside the domain of $f_i$.
If $x$ and $y$ both lie in this support, then
\begin{align*}
 |p_i(x)-p_i(y)|
 &\le w_i(x)|f_i(x)-f_i(y)|
   +|w_i(x)-w_i(y)|\,|f_i(y)|\\
 &\le (rL+B)d(x,y).
\end{align*}
If neither lies in the support, both products vanish.  If only $x$
does, the reverse triangle inequality gives
$w_i(x)=r-d(x,c_i)\le d(x,y)$ because $d(y,c_i)\ge r$.
Thus $|p_i(x)|\le B\,d(x,y)$, which is also bounded by
$(rL+B)d(x,y)$.  The case with only $y$ is symmetric.  Therefore every
$p_i$ extends continuously across the support boundary and has the
stated global Lipschitz bound.

Put $P=\sum_i p_i$.  We have
\[
 |P(x)-P(y)|\le n(rL+B)d(x,y),\qquad |P(y)|\le B W(y).
\]
Use $W(x)\ge r/2$ and $|W(x)-W(y)|\le n\,d(x,y)$ to estimate
\begin{align*}
 \left|\frac{P(x)}{W(x)}-\frac{P(y)}{W(y)}\right|
 &\le \frac{|P(x)-P(y)|}{W(x)}
   +\frac{|P(y)|\,|W(y)-W(x)|}{W(x)W(y)}\\
 &\le\frac{2n(rL+2B)}{r}d(x,y),
\end{align*}
which is \eqref{eq:gluing-lip}.  These rational-query estimates give
unique extensions to $X$ and remain true there by passage along
regular names with the displayed modulus.  The range bound follows
from $|P(x)|\le BW(x)$.  Finally, at a rational query with
$d(x,c_j)<r$, every index contributing
to the weighted average also has $d(x,c_i)<r$.  The overlap hypothesis
therefore gives
\[
 |g(x)-f_j(x)|
 \le\sum_i\lambda_i(x)|f_i(x)-f_j(x)|\le\alpha.
\]
For an arbitrary $x\in X$ with $d(x,c_j)<r$, sufficiently accurate
rational approximations lie in the same guarded ball; the Lipschitz
bounds for $g$ and $f_j$ transfer this estimate to $x$.
Only positivity, the triangle inequality, and finite sums were used,
so the proof is unchanged for complex values.
\end{proof}

For an implementation on the rational basis, the support test compares
two rationals.  The uniform arithmetic theorem uses one fixed total
PR program $E(i,x,t)$ for the patch parameter $i<n$, whose correctness
is required only inside the corresponding guard.  An externally
fixed finite collection of local programs may instead be compiled into
such an evaluator.  This does not permit unbounded interpretation of
arbitrary input program codes.  The program never
uses a local value outside the support of its weight.  The larger
radius $2r$ supplies room for evaluating names near support boundaries;
the preceding global modulus then extends the rational construction to
$X$.  For real queries one uses this extension, not an exact decision
whether $d(x,c_i)=r$.

More explicitly, if $E_i(x,t)$ approximates $f_i(x)$ within $2^{-t}$
on its guard, evaluate
\[
 G(x,t)=\frac{\sum_{i:\,d(x,c_i)<r}w_i(x)E_i(x,t)}{W(x)}.
\]
Every rational query to a local evaluator is within its guard, and
$|G(x,t)-g(x)|\le2^{-t}$ because the nonnegative normalized weights
sum to one.  Thus the gluing construction also supplies an explicit
regular-name evaluator, not just a semantic existence assertion.

This solves a specific presentation problem: local continuous data
with a certified finite metric cover can be glued by a concrete
formula, and the reciprocal bound and continuity modulus are computed
from the cover.  The construction does not manufacture finite covers
from an arbitrary open-cover enumeration, or smooth local data from
mere continuous data.

\subsection{Why regular subsequence extraction remains unavailable}

The distinction between approximation and subsequence selection is
already visible for constant functions.  Here is a self-contained
computability obstruction of the type introduced by Specker
\cite{Specker1949}; it also explains why one should keep the different
Ascoli principles separate \cite{Kreuzer2014}.

Fix a bounded PR simulator for Turing machines and set
\[
 q_s=\sum_{\substack{e<s\\ \text{machine }e\text{ on }e
                      \text{ halts by step }s}}4^{-(e+1)}.
\]
This is a PR sequence of rational numbers, is increasing, and lies in
$[0,1/3]$.  In the external classical interpretation its limit is
\[
 q=\sum_{e\in K}4^{-(e+1)},
\]
where $K$ is the diagonal halting set.  This real is not computable.
Indeed, after the bits of $K$ below $e$ have been determined, subtract
their contributions.  If $e\notin K$, the remaining value is at most
$4^{-(e+1)}/3$, whereas if $e\in K$ it is at least $4^{-(e+1)}$.
An approximation with error less than $4^{-(e+1)}/6$ distinguishes
these cases by comparison with $2\,4^{-(e+1)}/3$.  A computable name
for $q$ would consequently decide $K$, a contradiction.

Now let $f_s(x)=q_s$ on $[0,1]$.  The family has bound $B=1$ and the
common modulus $\mu(j)=0$ for every $j$.  Its finite approximants are
particularly easy.  Nevertheless, a computable strictly increasing
subsequence together with a computable modulus of uniform convergence
would compute $q$: evaluate a constant subsequence member at the index
specified by the modulus.  Thus no such regular subsequence can be
produced in general.  This is an external impossibility theorem about
uniform computations, not a new internally formalized incompleteness
claim for \ipra.  It does not rule out a particular sequence having a
supplied modulus or a problem-specific compactness argument producing
one.

\subsection{The exact rational target and a Lean development path}

For implementation, the semantic functions $f_a$ can be replaced
throughout by the following rational contracts, where $x,y\in A$ and
$p,q,j\in\N$:
\begin{align}
 |E(a,x,p)-E(a,x,q)|&\le d_p+d_q,\label{eq:top-name}\\
 |E(a,x,p)|&\le B+d_p,\label{eq:top-range}\\
 d(x,y)\le2^{-\mu(j)}\quad&\Longrightarrow\quad
 |E(a,x,p)-E(a,y,q)|\le d_j+d_p+d_q.
 \label{eq:top-space}
\end{align}
Together with the metric and net laws, these are sufficient for the
finite approximation theorem.  They state precisely what a rational
evaluation interface must provide.  They are hypotheses about fixed
programs, not automatically true properties of arbitrary program
codes.

At precision $k$, use $e_i=E(a,c(s_k,i),k+3)$ and the clipped floor
defined above.  Since \eqref{eq:top-range} gives
$e_i\in[-B-\delta_k,B+\delta_k]$, a direct three-case comparison
shows
\[
 |h_i-e_i|\le\delta_k.
\]
For $e_i\in[-B,B]$ this is the rounding bound.  For $e_i>B$ or
$e_i<-B$, clipping returns the respective endpoint and the range
contract gives the same bound.  If $w_i(x)>0$,
\eqref{eq:top-space} with $j=k+2$ now gives
\[
 |h_i-E(a,x,p)|
 \le\delta_k+d_{k+2}+d_{k+3}+d_p
 =d_{k+1}+d_p.
\]
Averaging proves the exact first-order target
\begin{equation}\label{eq:top-object-target}
 \boxed{\quad
 |g_k(a,x)-E(a,x,p)|\le 2^{-(k+1)}+2^{-p}.
 \quad}
\end{equation}
Every term in this formula is a rational code.  After checking
positive denominators, multiplication clears them, yielding a
quantifier-free integer inequality.  The net witness needed for the
denominator is an existential premise used outside the induction over
the finite sum.  In particular there is no induction on a statement
quantifying over all functions, all real points, or all precisions.

A Lean implementation should use two deliberately separate layers.
The following list gives the definitions and theorem order; the names
are proposed module names, not claims about compiled code delivered
elsewhere in this repository.
\begin{enumerate}
 \item \texttt{RationalMetricBasis} contains a decidable domain, a
 rational distance program, the pseudometric laws, \texttt{netSize},
 \texttt{netPoint}, and \eqref{eq:top-net}.  The arithmetic version
 stores PR syntax and object derivations for those laws; an ordinary
 Lean record with rational functions is the semantic specification.
 \item \texttt{TentWeights} defines the rational positive part,
 weights, denominator, and coefficient reconstruction.  Prove
 \texttt{weight\_nonnegative}, \texttt{denominator\_lower},
 \texttt{weight\_difference}, \texttt{partition\_variation}, and
 \texttt{reconstruction\_nonexpansive}, in that order.  Finite-sum
 estimates use existing sum linearity and rational-order derivations.
 \item \texttt{EquicontinuousEvaluator} stores $B,\mu,E$ and the
 three contracts \eqref{eq:top-name}--\eqref{eq:top-space}.  Add
 \texttt{clipFloor}, its grid-membership and error lemmas, and
 mixed-radix encoding of coefficient lists.
 \item \texttt{AscoliApproximation} defines $s_k,r_k,n_k,M_k,H_k$
 and $g_k$ literally as above.  Derive
 \texttt{ascoli\_finite\_error} with target
 \eqref{eq:top-object-target}, and separately the grid-code bound.
 No existential over function names is needed: $g_k$ is an explicit
 rational expression from the given observations.
 \item \texttt{GuardedGluing} handles the two support cases for each
 product, proves its Lipschitz bound, then sums and divides using
 \texttt{denominator\_lower}.  A semantic completion theorem
 identifies the resulting regular name with the continuous gluing.
 \item \texttt{AscoliConverse} introduces a bounded covering index,
 uses the three-term triangle estimate, and removes the index by
 existential elimination.  The maximum of a finite list supplies
 \eqref{eq:ascoli-converse-modulus}.
\end{enumerate}

For the object endpoint, the goal is an actual derivation
$\Gamma\vdash[\mathrm{PRAAx}]\varphi$ after inserting the chosen
PR metric, net, and evaluator terms.  The context $\Gamma$ retains
their explicit contracts.  A semantic Lean proof about
\texttt{MetricSpace} or \texttt{ContinuousMap} is useful for the
external interpretation but cannot by itself discharge this endpoint.
The required arithmetic induction invariants concern finite partial
sums and bounded coefficient lists and therefore fit the existing
quantifier-free derivation discipline.
